\subsection{Linear Operators as Modules}

Throughout this subsection, $\mathcal{V}$ is finite-dimensional over $\mathbb{F}$ unless explicitly stated otherwise.

\begin{perspective} \label{per:operator-as-polynomial-action} \hlt{An Operator as a Polynomial Action}\\
Fix $T\in\LL(\mathcal{V})$. Evaluation at $T$ defines an $\mathbb{F}[z]$-action on $\mathcal{V}$ by $p\cdot\mathbf{v}\defeq p(T)\mathbf{v}$.\\
Thus the orbit $\mathbf{v},T\mathbf{v},T^2\mathbf{v},\ldots$ is encoded by the cyclic subspace $\{p(T)\mathbf{v}:p\in\mathbb{F}[z]\}=\operatorname{span}\{T^j\mathbf{v}:j\geq0\}$.
\end{perspective}

\begin{definition} \label{def:operator-module-and-matrix-module} \hlt{The Module Attached to an Operator}\\
Let $T\in\LL(\mathcal{V})$. The \hlt{operator module} attached to $T$ is the $\mathbb{F}[z]$-module $\mathcal{V}_T$ with underlying additive group $\mathcal{V}$ and action $p(z)\cdot\mathbf{v}\defeq p(T)\mathbf{v}$.
For $A\in\operatorname{Mat}(n,n,\mathbb{F})$, define $(\mathbb{F}^n)_A$ analogously by $p(z)\cdot\mathbf{x}\defeq p(A)\mathbf{x}$.
\end{definition}

\begin{proposition} \label{prop:operator-modules-are-fg-torsion} \hlt{Operator Modules Are Finitely Generated Torsion Modules}\\
Let $0<\dim\mathcal{V}<\infty$ and $T\in\LL(\mathcal{V})$. Then $\mathcal{V}_T$ is a finitely generated torsion $\mathbb{F}[z]$-module.
\begin{proof}
Let $(\mathbf{v}_1,\ldots,\mathbf{v}_n)$ be an $\mathbb{F}$-basis of $\mathcal{V}$. Since $\mathbb{F}\subseteq\mathbb{F}[z]$, this basis generates $\mathcal{V}_T$.\\
By Proposition \ref{prop:finite-dimensional-operators-have-annihilating-polynomials}, there exists $0\neq p\in\mathbb{F}[z]$ with $p(T)=0$. Thus $p\cdot\mathbf{v}=p(T)\mathbf{v}=\mathbf{0}$ for all $\mathbf{v}\in\mathcal{V}$, so $\mathcal{V}_T$ is torsion. Finally, $\mathbb{F}[z]$ is a PID by Theorem \ref{thm:integers-and-polynomial-ring-are-pids}.\\
\hlt{Structural Commentary}: Finite-dimensionality gives a global polynomial relation for the operator.\\
\hlt{Key Insight}: Rational canonical form is the coordinate form of the structure theorem for $\mathcal{V}_T$.
\end{proof}
\end{proposition}

\begin{theorem} \label{thm:operator-submodules-invariant-subspaces} \hlt{Submodules Are Invariant Subspaces}\\
Let $T\in\LL(\mathcal{V})$ and $\mathcal{S}\subseteq\mathcal{V}$. Then $\mathcal{S}\leq_{\mathbb{F}[z]}\mathcal{V}_T$ if and only if $\mathcal{S}\leq_{\mathbb{F}}\mathcal{V}$ and $T\mathcal{S}\subseteq\mathcal{S}$.
\begin{proof}
\begin{enumerate}[label=\roman*.]
\setlength{\itemsep}{0pt}
\item[$(\Rightarrow)$] If $\mathcal{S}\leq_{\mathbb{F}[z]}\mathcal{V}_T$, then $\mathbb{F}\subseteq\mathbb{F}[z]$ gives $\mathcal{S}\leq_{\mathbb{F}}\mathcal{V}$. Since $z\cdot\mathbf{s}=T\mathbf{s}$ for $\mathbf{s}\in\mathcal{S}$, we have $T\mathcal{S}\subseteq\mathcal{S}$.
\item[$(\Leftarrow)$] If $\mathcal{S}\leq_{\mathbb{F}}\mathcal{V}$ and $T\mathcal{S}\subseteq\mathcal{S}$, then $T^j\mathcal{S}\subseteq\mathcal{S}$ for all $j\geq0$.\\
Hence for $p(z)=\sum_ja_jz^j$ and $\mathbf{s}\in\mathcal{S}$, $p\cdot\mathbf{s}=p(T)\mathbf{s}=\sum_ja_jT^j\mathbf{s}\in\mathcal{S}$, so $\mathcal{S}\leq_{\mathbb{F}[z]}\mathcal{V}_T$.
\end{enumerate}
\hlt{Structural Commentary}: The single scalar $z$ remembers the operator.\\
\hlt{Key Insight}: Invariance under $T$ is exactly closure under all polynomial scalars.
\end{proof}
\end{theorem}

\begin{theorem} \label{thm:operator-similarity-iff-module-isomorphism} \hlt{Similarity Is Module Isomorphism}\\
Let $T,S\in\LL(\mathcal{V})$. Then $\mathcal{V}_T\cong_{\mathbb{F}[z]}\mathcal{V}_S\Longleftrightarrow T\sim S$. More precisely, for $\Phi\in\operatorname{GL}(\mathcal{V})$, $\Phi:\mathcal{V}_T\to\mathcal{V}_S$ is $\mathbb{F}[z]$-linear if and only if $\Phi T=S\Phi$, equivalently $S=\Phi T\Phi^{-1}$.
\begin{proof}
\begin{enumerate}[label=\roman*.]
\setlength{\itemsep}{0pt}
\item[$(\Rightarrow)$] If $\Phi:\mathcal{V}_T\to\mathcal{V}_S$ is an $\mathbb{F}[z]$-isomorphism, then $\Phi\in\operatorname{GL}(\mathcal{V})$, $\Phi(T\mathbf{v})=\Phi(z\cdot\mathbf{v})=z\cdot\Phi(\mathbf{v})=S\Phi(\mathbf{v})$ $\forall \mathbf{v}\in\mathcal{V}$.\\
Hence $\Phi T=S\Phi$, so $S=\Phi T\Phi^{-1}$.
\item[$(\Leftarrow)$] If $\Phi\in\operatorname{GL}(\mathcal{V})$ and $S=\Phi T\Phi^{-1}$, then $\Phi T=S\Phi$, hence $\Phi T^j=S^j\Phi$ for all $j\geq0$.\\
Thus for $p(z)=\sum_ja_jz^j$, $\Phi(p(T)\mathbf{v})=p(S)\Phi(\mathbf{v})$, so $\Phi:\mathcal{V}_T\to\mathcal{V}_S$ is an $\mathbb{F}[z]$-isomorphism.
\end{enumerate}
\hlt{Structural Commentary}: A module isomorphism is precisely an intertwiner for the action of $z$.\\
\hlt{Key Insight}: Similarity is isomorphism of the associated $\mathbb{F}[z]$-modules.
\end{proof}
\end{theorem}

\begin{corollary} \label{cor:matrix-similarity-iff-matrix-modules-isomorphic} \hlt{Matrix Similarity Is Matrix-Module Isomorphism}\\
Let $A,B\in\operatorname{Mat}(n,n,\mathbb{F})$. Then $A\sim B\Longleftrightarrow(\mathbb{F}^n)_A\cong_{\mathbb{F}[z]}(\mathbb{F}^n)_B$. More precisely, for $X\in\operatorname{GL}_n(\mathbb{F})$, $X:(\mathbb{F}^n)_A\to(\mathbb{F}^n)_B$ is $\mathbb{F}[z]$-linear if and only if $XA=BX$, equivalently $B=XAX^{-1}$.
\begin{proof}
\begin{enumerate}[label=\roman*.]
\setlength{\itemsep}{0pt}
\item[$(\Rightarrow)$] If $X:(\mathbb{F}^n)_A\to(\mathbb{F}^n)_B$ is an $\mathbb{F}[z]$-isomorphism, then $X\in\operatorname{GL}_n(\mathbb{F})$ and $X(A\mathbf{x})=X(z\cdot\mathbf{x})=z\cdot X\mathbf{x}=BX\mathbf{x}$ for all $\mathbf{x}\in\mathbb{F}^n$. Hence $XA=BX$, so $B=XAX^{-1}$.
\item[$(\Leftarrow)$] If $X\in\operatorname{GL}_n(\mathbb{F})$ and $B=XAX^{-1}$, then $XA=BX$, hence $XA^j=B^jX$ for all $j\geq0$. Thus $Xp(A)=p(B)X$ for all $p\in\mathbb{F}[z]$, so $X:(\mathbb{F}^n)_A\to(\mathbb{F}^n)_B$ is an $\mathbb{F}[z]$-isomorphism.
\end{enumerate}
\hlt{Structural Commentary}: Similarity is invertible module intertwining.\\
\hlt{Key Insight}: The matrix is remembered by the scalar action of $z$.
\end{proof}
\end{corollary}

\begin{definition} \label{def:operator-module-annihilator-minimal-polynomial} \hlt{Module Annihilator and Minimal Polynomial}
\begin{enumerate}[label=(\roman*)]
\setlength{\itemsep}{0pt}
\item The \hlt{annihilator} of $\mathcal{V}_T$ is $\operatorname{Ann}_{\mathbb{F}[z]}(\mathcal{V}_T)\defeq\{p\in\mathbb{F}[z]:p(T)\mathcal{V}=\{\mathbf{0}\}\}$.
\item \hlt{Minimal polynomial} $\mu_T$ is unique monic generator of this nonzero principal ideal, so $\operatorname{Ann}_{\mathbb{F}[z]}(\mathcal{V}_T)=\langle\mu_T\rangle$.
\item For $A\in\operatorname{Mat}(n,n,\mathbb{F})$, define $\mu_A\defeq\mu_{T_A}$, where $T_A(\mathbf{x})=A\mathbf{x}$.
\end{enumerate}
\end{definition}

\begin{proposition} \label{prop:minimal-polynomial-module-order} \hlt{Minimal Polynomial as Module Order}\\
Let $T\in\LL(\mathcal{V})$. Then $\mu_T$ is the unique monic polynomial of minimal degree in $\{p\in\mathbb{F}[z]:p(T)=0\}$. If $S\sim T$, then $\mu_S=\mu_T$; hence similar matrices have the same minimal polynomial.
\begin{proof}
By Definition \ref{def:operator-module-annihilator-minimal-polynomial}, $\operatorname{Ann}_{\mathbb{F}[z]}(\mathcal{V}_T)=\langle\mu_T\rangle$ and $p\in\operatorname{Ann}_{\mathbb{F}[z]}(\mathcal{V}_T)$ iff $p(T)=0$. Thus $\mu_T$ is the unique monic annihilator of minimal degree. If $S=\Phi T\Phi^{-1}$, then $p(S)=\Phi p(T)\Phi^{-1}$ for all $p\in\mathbb{F}[z]$, so $\operatorname{Ann}_{\mathbb{F}[z]}(\mathcal{V}_S)=\operatorname{Ann}_{\mathbb{F}[z]}(\mathcal{V}_T)$ and $\mu_S=\mu_T$. The matrix case follows by applying this to the associated operators.\\
\hlt{Structural Commentary}: The minimal polynomial is the order of the whole operator module.\\
\hlt{Key Insight}: Similarity preserves the module and therefore preserves its annihilator ideal.
\end{proof}
\end{proposition}

\begin{definition} \label{def:cyclic-subspace-and-vector-annihilator} \hlt{Cyclic Subspace and Vector Annihilator}\\
Let $T\in\LL(\mathcal{V})$ and $\mathbf{v}\in\mathcal{V}$.
\begin{enumerate}[label=(\roman*)]
\setlength{\itemsep}{0pt}
\item The \hlt{$T$-cyclic subspace generated by $\mathbf{v}$} is $\mathcal{C}_T(\mathbf{v})\defeq\operatorname{span}\{T^j\mathbf{v}:j\geq0\}$
\item The vector $\mathbf{v}$ is a \hlt{cyclic vector} for $T$ if $\mathcal{C}_T(\mathbf{v})=\mathcal{V}$.
\item The \hlt{annihilator of $\mathbf{v}$ for $T$} is $\operatorname{ann}_T(\mathbf{v})\defeq\{p\in\mathbb{F}[z]:p(T)\mathbf{v}=\mathbf{0}\}$
\end{enumerate}
\end{definition}

\begin{proposition} \label{prop:cyclic-subspace-smallest-invariant-subspace} \hlt{Cyclic Subspaces Are Smallest Invariant Hosts}\\
Let $T\in\LL(\mathcal{V})$ and $\mathbf{v}\in\mathcal{V}$. Then $\mathcal{C}_T(\mathbf{v})=\{p(T)\mathbf{v}:p\in\mathbb{F}[z]\}$.\\
Moreover, $\mathcal{C}_T(\mathbf{v})$ is the smallest $T$-invariant subspace containing $\mathbf{v}$, and $\operatorname{ann}_T(\mathbf{v})$ is an ideal of $\mathbb{F}[z]$.
\begin{proof} Equality follows as $p(T)\mathbf{v}$ is a finite linear combination of the orbit vectors $T^j\mathbf{v}$, and every such finite linear combination equals $p(T)\mathbf{v}$ for some $p\in\mathbb{F}[z]$. Also $T(p(T)\mathbf{v})=(zp)(T)\mathbf{v}$, so $\mathcal{C}_T(\mathbf{v})$ is $T$-invariant.\\
If $\mathcal{U}$ is $T$-invariant and contains $\mathbf{v}$, then $T^j\mathbf{v}\in\mathcal{U}$ for every $j\geq0$, hence $\mathcal{C}_T(\mathbf{v})\subseteq\mathcal{U}$. Finally, if $p,q\in\operatorname{ann}_T(\mathbf{v})$ and $r\in\mathbb{F}[z]$, then $(p+q)(T)\mathbf{v}=\mathbf{0}$ and $(rp)(T)\mathbf{v}=r(T)p(T)\mathbf{v}=\mathbf{0}$ by Proposition \ref{prop:operator-polynomial-products}. $\operatorname{ann}_T(\mathbf{v})$ is an ideal.\\
\hlt{Structural Commentary}: A cyclic subspace is the span of the orbit generated by repeatedly applying $T$.\\
\hlt{Key Insight}: Invariance is exactly what lets the whole orbit stay inside the same host subspace.
\end{proof}
\end{proposition}

\begin{figure}[H]
  \centering
  \begin{tikzpicture}[
  >=Stealth,
  every node/.style={font=\small},
  orbit/.style={draw=blue!60!black, fill=blue!6, rounded corners=2pt,
    minimum width=1.15cm, minimum height=0.58cm, align=center, inner sep=4pt},
  subspace/.style={draw=dartmouthgreen!65!black, fill=dartmouthgreen!8, rounded corners=2pt,
    minimum width=4.7cm, minimum height=0.72cm, align=center, inner sep=6pt},
  arr/.style={->, line width=0.7pt, draw=black!70}
]
  \node[orbit] (v0) at (-2.45,0) {$\mathbf{v}$};
  \node[orbit] (v1) at (-0.85,0) {$T\mathbf{v}$};
  \node[orbit] (v2) at (0.85,0) {$T^2\mathbf{v}$};
  \node[orbit, minimum width=0.8cm] (dots) at (2.35,0) {$\cdots$};

  \draw[arr] (v0) -- node[above] {$T$} (v1);
  \draw[arr] (v1) -- node[above] {$T$} (v2);
  \draw[arr] (v2) -- node[above] {$T$} (dots);

  \node[subspace] (C) at (0,-1.35)
    {$\mathcal{C}_T(\mathbf{v})=\operatorname{span}\{T^j\mathbf{v}:j\geq0\}$};

  \draw[arr] (0,-0.38) -- node[right] {$\operatorname{span}$} (C.north);
\end{tikzpicture}

  \caption{The $T$-orbit of $\mathbf{v}$ spans $\mathcal{C}_T(\mathbf{v})$.}
  \label{fig:operator-cyclic-subspace-polynomial-action}
\end{figure}

\begin{definition} \label{def:vector-order-under-operator} \hlt{Order of a Vector under an Operator}\\
Let $\mathcal{V}$ be finite-dimensional, let $T\in\LL(\mathcal{V})$, and let $\mathbf{v}\neq\mathbf{0}$. The \hlt{order of $\mathbf{v}$ under $T$} is the unique monic polynomial $m_{\mathbf{v}}\in\operatorname{ann}_T(\mathbf{v})$ of least degree; equivalently, $\operatorname{ann}_T(\mathbf{v})=(m_{\mathbf{v}})$
\end{definition}

\begin{definition} \label{def:companion-matrix} \hlt{Companion Matrix}\\
Let $f(z)=z^n+c_{n-1}z^{n-1}+\cdots+c_1z+c_0$ be monic. For $n\geq2$, the \hlt{companion matrix} of $f$ is
\begin{align}
C_f\defeq
\begin{bmatrix}
0&0&\cdots&0&-c_0\\
1&0&\cdots&0&-c_1\\
0&1&\cdots&0&-c_2\\
\vdots&\vdots&\ddots&\vdots&\vdots\\
0&0&\cdots&1&-c_{n-1}
\end{bmatrix}.
\nonumber
\end{align}
For $n=1$, set $C_f=[-c_0]$.
\end{definition}

\begin{theorem} \label{thm:companion-matrix-has-prescribed-minimal-polynomial} \hlt{Companion Matrix Realises a Polynomial}\\
Let $f(z)=z^n+\sum_{i=0}^{n-1}c_iz^i$ be monic and let $C_f$ be its companion matrix. Then $p_{C_f}=\mu_{C_f}=f$.
\begin{proof}
Let $(\mathbf{e}_1,\ldots,\mathbf{e}_n)$ be the standard basis.\\
For $0\leq j<n$, $C_f^j\mathbf{e}_1=\mathbf{e}_{j+1}$, and last column gives $C_f^n\mathbf{e}_1=-\sum\limits_{i=0}^{n-1}c_iC_f^i\mathbf{e}_1$. Hence $f(C_f)\mathbf{e}_1=0$.\\
Since $\mathbf{e}_j=C_f^{j-1}\mathbf{e}_1$ and $f(C_f)C_f=C_ff(C_f)$, $f(C_f)\mathbf{e}_j=0$ $\forall j$, so $f(C_f)=0$.\\
If $0\neq g(z)=\sum\limits_{i=0}^m b_iz^i$ with $m<n$, then $g(C_f)\mathbf{e}_1=\sum\limits_{i=0}^m b_i\mathbf{e}_{i+1}\neq\mathbf{0}$.\\
Thus no nonzero polynomial of degree $<n$ annihilates $C_f$, so $\deg\mu_{C_f}=n$.\\
Since $\mu_{C_f}\mid f$ and both are monic of degree $n$, $\mu_{C_f}=f$.\\
By Cayley--Hamilton, $\mu_{C_f}\mid p_{C_f}$; as $p_{C_f}$ and $f$ are monic of degree $n$, $p_{C_f}=f$.\\
\hlt{Structural Commentary}: The cyclic vector $\mathbf{e}_1$ generates the basis, and last column imposes relation $f(C_f)=0$.\\
\hlt{Key Insight}: The companion matrix makes $f$ the first possible annihilating relation.
\end{proof}
\end{theorem}

\begin{theorem} \label{thm:cyclic-subspace-quotient-and-companion} \label{prop:cyclic-basis-gives-companion-matrix} \hlt{Cyclic Subspaces as Polynomial Quotients}\\
Let $\dim\mathcal{V}<\infty$, $T\in\LL(\mathcal{V})$, $0\neq\mathbf{v}\in\mathcal{V}$, and let $m_{\mathbf{v}}(z)=z^d+\sum\limits_{i=0}^{d-1}c_iz^i$ be the order of $\mathbf{v}$.\\
Define $\Phi_{\mathbf{v}}:\mathbb{F}[z]\to\mathcal{C}_T(\mathbf{v})$ by $\Phi_{\mathbf{v}}(p)=p(T)\mathbf{v}$ and $\mathcal{B}_{\mathbf{v}}\defeq(\mathbf{v},T\mathbf{v},\ldots,T^{d-1}\mathbf{v})$. Then
\begin{enumerate}[label=(\roman*)]
\setlength{\itemsep}{0pt}
\item $\operatorname{im}\Phi_{\mathbf{v}}=\mathcal{C}_T(\mathbf{v})$, $\ker\Phi_{\mathbf{v}}=\langle m_{\mathbf{v}}\rangle$, and $\mathcal{C}_T(\mathbf{v})\cong\mathbb{F}[z]/\langle m_{\mathbf{v}}\rangle$;
\item $\mathcal{B}_{\mathbf{v}}$ is a basis of $\mathcal{C}_T(\mathbf{v})$;
\item $[T|_{\mathcal{C}_T(\mathbf{v})}]_{\mathcal{B}_{\mathbf{v}}}=C_{m_{\mathbf{v}}}$.
\end{enumerate}
\begin{proof}
\begin{enumerate}[label=(\roman*)]
\setlength{\itemsep}{0pt}
\item By Proposition \ref{prop:cyclic-subspace-smallest-invariant-subspace}, $\operatorname{im}\Phi_{\mathbf{v}}=\mathcal{C}_T(\mathbf{v})$. $\ker\Phi_{\mathbf{v}}=\operatorname{ann}_T(\mathbf{v})=\langle m_{\mathbf{v}}\rangle$, by Theorem \ref{thm:quotient-by-kernel-is-range} $\mathcal{C}_T(\mathbf{v})\cong\mathbb{F}[z]/\langle m_{\mathbf{v}}\rangle$.
\item If $\sum\limits_{j=0}^{d-1}a_jT^j\mathbf{v}=\mathbf{0}$, then $q(z)=\sum\limits_{j=0}^{d-1}a_jz^j\in\operatorname{ann}_T(\mathbf{v})$ and $\deg q<d$, hence $q=0$; thus $\mathcal{B}_{\mathbf{v}}$ is independent.\\
For $p=hm_{\mathbf{v}}+r$ with $\deg r<d$ or $r=0$, $p(T)\mathbf{v}=r(T)\mathbf{v}\in\operatorname{span}\mathcal{B}_{\mathbf{v}}$, so $\mathcal{B}_{\mathbf{v}}$ spans.
\item For $0\leq j<d-1$, $T(T^j\mathbf{v})=T^{j+1}\mathbf{v}$. Since $m_{\mathbf{v}}(T)\mathbf{v}=\mathbf{0}$, $T^d\mathbf{v}=-\sum\limits_{i=0}^{d-1}c_iT^i\mathbf{v}$, last-column relation of $C_{m_{\mathbf{v}}}$.
\end{enumerate}
\hlt{Structural Commentary}: The first relation cuts out the quotient and companion block.\\
\hlt{Key Insight}: One orbit carries the whole cyclic subspace.
\end{proof}
\end{theorem}

\begin{perspective} \label{per:module-operator-dictionary} \hlt{The Translation Dictionary}\\
The module vocabulary for $\mathcal{V}_T$ and the vector-space vocabulary for $T$ match as follows:
\begin{center}
\begin{tabular}{c|c}
$\mathbb{F}[z]$-module $\mathcal{V}_T$ & $\mathbb{F}$-vector space $\mathcal{V}$\\
\hline
$p(z)\cdot\mathbf{v}$ & $p(T)\mathbf{v}$\\
submodule & $T$-invariant subspace\\
$\operatorname{Ann}_{\mathbb{F}[z]}(\mathcal{V}_T)=\langle\mu_T\rangle$ & minimal polynomial of $T$\\
$\langle\!\langle\mathbf{v}\rangle\!\rangle_{\mathbb{F}[z]}$ & $\mathcal{C}_T(\mathbf{v})$\\
$\mathbb{F}[z]/\langle m_{\mathbf{v}}\rangle$ & cyclic block with companion matrix $C_{m_{\mathbf{v}}}$
\end{tabular}
\end{center}
\hlt{Key Insight}: Once $z$ is interpreted as $T$, module structure becomes operator structure.
\end{perspective}

\begin{method} \label{meth:operator-module-to-spectral-form} \hlt{Moving Between Module and Spectral Data}\\
For $T\in\LL(\mathcal{V})$, translate the $\mathbb{F}[z]$-module data of $\mathcal{V}_T$ into spectral data as follows.
\begin{enumerate}[label=(\roman*)]
\setlength{\itemsep}{0pt}
\item Factor $\mu_T=\prod_{i=1}^s p_i^{e_i}$ over $\mathbb{F}$; the primary components are $\ker p_i(T)^{e_i}$.
\item Decompose each $\ker p_i(T)^{e_i}$ into cyclic summands; their orders $p_i^{e_{i,j}}$ are the elementary divisors.
\item Over $\mathbb{F}$, each elementary divisor $p_i^{e_{i,j}}$ gives a companion block $C_{p_i^{e_{i,j}}}$ in rational canonical form.
\item Over splitting field for $p_i$, same primary cyclic summands refine into Jordan blocks (Perspective \ref{per:elementary-divisors-and-jordan-blocks}).
\item $T$ is diagonalizable exactly when every elementary divisor is linear, i.e., Corollary \ref{cor:diagonal-rational-canonical-form-criterion} applies.
\end{enumerate}
\hlt{Key Insight}: Rational and Jordan forms encode the same elementary-divisor data over different fields.
\end{method}

\begin{definition} \label{def:operator-invariant-factors-elementary-divisors} \hlt{Invariant Factors and Elementary Divisors}\\
Let $T\in\LL(\mathcal{V})$.
\begin{enumerate}[label=(\roman*)]
\setlength{\itemsep}{0pt}
\item A \hlt{cyclic decomposition} is $\mathcal{V}=\bigoplus_{j=1}^r\mathcal{C}_j$, where each $\mathcal{C}_j$ is nonzero, $T$-invariant, and cyclic.
\item $\operatorname{InvFact}(T)\defeq(f_1,\ldots,f_r)$, where $f_j$ are monic nonconstant invariant factors of $\mathcal{V}_T$; $f_r\mid f_{r-1}\mid\cdots\mid f_1$.
\item $\operatorname{ElemDiv}(T)$ is the multiset of monic prime powers in the $f_j$, counted with multiplicity.
\item \hlt{Elementary divisor basis} is $\mathcal{B}_{\operatorname{ED}}\defeq\bigsqcup_{\alpha}(\mathbf{v}_{\alpha},T\mathbf{v}_{\alpha},\ldots,T^{d_{\alpha}-1}\mathbf{v}_{\alpha})$, where $d_{\alpha}=\deg m_{\mathbf{v}_{\alpha}}$; $\{m_{\mathbf{v}_{\alpha}}\}_{\alpha}=\operatorname{ElemDiv}(T)$.
\item For $A\in\operatorname{Mat}(n,n,\mathbb{F})$, set $\operatorname{InvFact}(A)\defeq\operatorname{InvFact}(T_A)$ and $\operatorname{ElemDiv}(A)\defeq\operatorname{ElemDiv}(T_A)$.
\end{enumerate}
\end{definition}

\begin{theorem} \label{thm:operator-primary-cyclic-decomposition} \hlt{Primary Cyclic Decomposition for an Operator}\\
Let $\mu_T=\prod_{i=1}^sp_i^{e_i}$ with $p_i\in\mathbb{F}[z]$ distinct monic irreducible, and set $\mathcal{V}_{p_i}\defeq\ker p_i(T)^{e_i}$. Then
\begin{enumerate}[label=(\roman*)]
\setlength{\itemsep}{0pt}
\item $\mathcal{V}_T=\bigoplus_{i=1}^s\mathcal{V}_{p_i}$;
\item $\mathcal{V}_{p_i}=\bigoplus_{j=1}^{k_i}\mathcal{C}_T(\mathbf{v}_{i,j})$, with $m_{\mathbf{v}_{i,j}}=p_i^{e_{i,j}}$ and $e_i=e_{i,1}\geq e_{i,2}\geq\cdots\geq e_{i,k_i}$;
\item $\operatorname{ElemDiv}(T)=\{p_i^{e_{i,j}}:1\leq i\leq s,\ 1\leq j\leq k_i\}$, corresponding cyclic bases form elementary divisor basis.
\end{enumerate}
\begin{proof}
\begin{enumerate}[label=(\roman*)]
\setlength{\itemsep}{0pt}
\item By Theorem \ref{thm:primary-cyclic-decomposition-pid}, $\mathcal{V}_T$, $\mathcal{V}_T=\bigoplus_i\mathcal{V}_T[p_i^\infty]$; $\mathcal{V}_T[p_i^\infty]\defeq\{\mathbf{x}\in\mathcal{V}:\exists N\geq1,\ p_i(T)^N\mathbf{x}=0\}=\ker p_i(T)^{e_i}=\mathcal{V}_{p_i}$.
\item By Theorem \ref{thm:primary-cyclic-decomposition-pid}, $\mathcal{V}_{p_i}=\bigoplus_{j=1}^{k_i}\mathcal{C}_T(\mathbf{v}_{i,j})$, $m_{\mathbf{v}_{i,j}}=p_i^{e_{i,j}}$ and $e_i=e_{i,1}\geq\cdots\geq e_{i,k_i}$.
\item Definition \ref{def:operator-invariant-factors-elementary-divisors} gives $\operatorname{ElemDiv}(T)=\{p_i^{e_{i,j}}:1\leq i\leq s,\ 1\leq j\leq k_i\}$, and concatenating the bases of $\mathcal{C}_T(\mathbf{v}_{i,j})$ gives an elementary divisor basis.
\end{enumerate}
\hlt{Structural Commentary}: Primary module pieces are polynomial-kernel pieces for the operator.\\
\hlt{Key Insight}: Elementary divisors are exactly the orders of the cyclic primary summands.
\end{proof}
\end{theorem}

\begin{theorem} \label{thm:operator-invariant-factor-skeleton} \hlt{Invariant-Factor Skeleton for Operators}\\
Let $0<\dim\mathcal{V}<\infty$, $T\in\LL(\mathcal{V})$, and $\operatorname{InvFact}(T)=(f_1,\ldots,f_r)$.\\
Then $f_r\mid\cdots\mid f_1$ and there are $T$-cyclic subspaces $\mathcal{C}_j$ such that $\mathcal{V}=\bigoplus_{j=1}^r\mathcal{C}_j$, $\mathcal{C}_j\cong\mathbb{F}[z]/\langle f_j\rangle$, $[T]=\bigoplus_{j=1}^r C_{f_j}$ in the concatenated cyclic bases, $\mu_T=f_1$, and $p_T=\prod_{j=1}^r f_j$.
\begin{proof}
By Proposition \ref{prop:operator-modules-are-fg-torsion}, $\mathcal{V}_T$ is finitely generated torsion over $\mathbb{F}[z]$.\\
Theorem \ref{thm:pid-invariant-factor-decomposition} gives $\mathcal{V}_T\cong\bigoplus_{j=1}^r\mathbb{F}[z]/\langle f_j\rangle$ with $f_r\mid\cdots\mid f_1$, uniquely.\\
Transport cyclic generators to $T$-cyclic subspaces $\mathcal{C}_j$; Theorem \ref{thm:cyclic-subspace-quotient-and-companion} gives $\mathcal{C}_j\cong\mathbb{F}[z]/\langle f_j\rangle$ and block $C_{f_j}$.\\
Since $\operatorname{Ann}(\mathcal{V}_T)=\langle\operatorname{lcm}(f_1,\ldots,f_r)\rangle=\langle f_1\rangle=\langle\mu_T\rangle$, $f_1=\mu_T$.\\
Finally, Theorem \ref{thm:companion-matrix-has-prescribed-minimal-polynomial} gives $p_{C_{f_j}}=f_j$, so $p_T=\prod_{j=1}^r f_j$.\\
\hlt{Structural Commentary}: Torsion module decomposition produces companion blocks.\\
\hlt{Key Insight}: Invariant factors are the polynomial-module version of choosing canonical blocks for similarity.
\end{proof}
\end{theorem}

\begin{theorem} \label{thm:reducing-decompositions-combine-elementary-divisors} \hlt{Reducing Decompositions Combine the Data}\\
Let $\mathcal{V}=\mathcal{A}\oplus\mathcal{B}$ with $\mathcal{A},\mathcal{B}$ $T$-invariant, and set $T_{\mathcal{A}}\defeq T|_{\mathcal{A}}$, $T_{\mathcal{B}}\defeq T|_{\mathcal{B}}$.\\
Then $\mu_T=\operatorname{lcm}(\mu_{T_{\mathcal{A}}},\mu_{T_{\mathcal{B}}})$ and $\operatorname{ElemDiv}(T)=\operatorname{ElemDiv}(T_{\mathcal{A}})\sqcup\operatorname{ElemDiv}(T_{\mathcal{B}})$, where $\sqcup$ is multiset union.\\
The invariant factors are obtained by recombining the elementary divisors into a divisibility chain.
\begin{proof}
By Theorem \ref{thm:operator-submodules-invariant-subspaces}, $\mathcal{V}_T=\mathcal{A}_{T_{\mathcal{A}}}\oplus\mathcal{B}_{T_{\mathcal{B}}}$ as $\mathbb{F}[z]$-modules.\\
Hence $\operatorname{Ann}(\mathcal{V}_T)=\operatorname{Ann}(\mathcal{A}_{T_{\mathcal{A}}})\cap\operatorname{Ann}(\mathcal{B}_{T_{\mathcal{B}}})=\langle\operatorname{lcm}(\mu_{T_{\mathcal{A}}},\mu_{T_{\mathcal{B}}})\rangle$.\\
Proposition \ref{prop:elementary-divisors-direct-sums} gives elementary-divisor identity.\\
\hlt{Structural Commentary}: Reducing decompositions split the operator module.\\
\hlt{Key Insight}: Direct sums take lcm on global orders and multiset union on elementary divisors.
\end{proof}
\end{theorem}

\begin{definition} \label{def:characteristic-polynomial-from-elementary-divisors} \hlt{Characteristic Polynomial from Elementary Divisors}\\
For $T\in\LL(\mathcal{V})$, define $\chi_T(z)\defeq\prod\limits_{q\in\operatorname{ElemDiv}(T)}q(z)$, with product over the multiset $\operatorname{ElemDiv}(T)$.\\
For $A\in\operatorname{Mat}(n,n,\mathbb{F})$, set $\chi_A\defeq\chi_{T_A}$.
\end{definition}

\begin{theorem} \label{thm:minimal-characteristic-elementary-divisor-relations} \hlt{Minimal and Characteristic Polynomials from Elementary Divisors}\\
Let $\operatorname{ElemDiv}(T)=\{p_i^{e_{i,j}}:1\leq i\leq s,\ 1\leq j\leq k_i\}$.\\
Then $\mu_T=\prod_i p_i^{\max_j e_{i,j}}$, $\chi_T=\prod_{i,j}p_i^{e_{i,j}}$, and $\deg\chi_T=\dim\mathcal{V}$.\\
Hence $\mu_T\mid\chi_T$, and $\mu_T,\chi_T$ have the same irreducible factors.
\begin{proof}
By Theorem \ref{thm:operator-primary-cyclic-decomposition}, $\mathcal{V}_T=\bigoplus_{i,j}\mathcal{C}_T(\mathbf{v}_{i,j})$ with $m_{\mathbf{v}_{i,j}}=p_i^{e_{i,j}}$.\\
Hence $\operatorname{Ann}(\mathcal{V}_T)=\langle\operatorname{lcm}_{i,j}p_i^{e_{i,j}}\rangle=\langle\prod_i p_i^{\max_j e_{i,j}}\rangle$, so $\mu_T=\prod_i p_i^{\max_j e_{i,j}}$. Definition \ref{def:characteristic-polynomial-from-elementary-divisors} gives $\chi_T=\prod_{i,j}p_i^{e_{i,j}}$.\\
By Theorem \ref{thm:cyclic-subspace-quotient-and-companion}, $\dim\mathcal{C}_T(\mathbf{v}_{i,j})=\deg p_i^{e_{i,j}}$; summing gives $\deg\chi_T=\dim\mathcal{V}$.\\
Divisibility and common irreducible factors are immediate.\\
\hlt{Structural Commentary}: $\mu_T$ records maximal exponents; $\chi_T$ records all exponents.\\
\hlt{Key Insight}: The pair $(\mu_T,\chi_T)$ forgets the block distribution.
\end{proof}
\end{theorem}

\begin{definition} \label{def:nonderogatory-operator-module} \hlt{Nonderogatory Operators}\\
$T\in\LL(\mathcal{V})$ is \hlt{nonderogatory} if $\mu_T=\chi_T$, equivalently $\deg\mu_T=\dim\mathcal{V}$.\\
For $A\in\operatorname{Mat}(n,n,\mathbb{F})$, $A$ is nonderogatory if $T_A$ is.
\end{definition}

\begin{theorem} \label{thm:nonderogatory-cyclic-equivalences} \hlt{Cyclic and Nonderogatory Are the Same}\\
For $T\in\LL(\mathcal{V})$, the following are equivalent:
\begin{enumerate}[label=(\roman*)]
\setlength{\itemsep}{0pt}
\item $\mathcal{V}_T$ is cyclic;
\item $T$ has a cyclic vector;
\item $\operatorname{ElemDiv}(T)=\{p_1^{e_1},\ldots,p_s^{e_s}\}$, one elementary divisor for each irreducible factor of $\mu_T$;
\item $T$ is nonderogatory.
\end{enumerate}
\begin{proof}
\begin{enumerate}[label=(\roman*).]
\setlength{\itemsep}{0pt}
\item[$(i)$]$\Rightarrow (ii)$ If $\mathcal{V}_T=\mathbb{F}[z]\mathbf{v}$, then $\mathcal{V}=\{p(T)\mathbf{v}:p\in\mathbb{F}[z]\}=\mathcal{C}_T(\mathbf{v})$, so $\mathbf{v}$ is cyclic.
\item[$(ii)$]$\Rightarrow (i)$ If $\mathcal{C}_T(\mathbf{v})=\mathcal{V}$, then $\mathcal{V}_T=\mathbb{F}[z]\mathbf{v}$.
\item[$(i)$]$\Leftrightarrow (iii)$ Theorem \ref{thm:cyclic-torsion-module-characterization} on $\mathcal{V}_T$. Cyclic torsion module has one cyclic summand for each primary factor.
\item[$(iii)$]$\Rightarrow (iv)$ If $\operatorname{ElemDiv}(T)=\{p_i^{e_i}\}_{i=1}^s$, then $\chi_T=\prod_i p_i^{e_i}=\mu_T$, so $T$ is nonderogatory.
\item[$(iv)$]$\Rightarrow (iii)$ Let $\operatorname{ElemDiv}(T)=\{p_i^{e_{i,j}}\}$. If $\mu_T=\chi_T$, then $\prod_i p_i^{\max_j e_{i,j}}=\prod_{i,j}p_i^{e_{i,j}}$, so $k_i=1$ for all $i$.
\end{enumerate}
\hlt{Structural Commentary}: Nonderogatory means all primary information can be generated by one vector.\\
\hlt{Key Insight}: Equality $\mu_T=\chi_T$ says that product and lcm of the cyclic orders coincide.
\end{proof}
\end{theorem}

\begin{perspective} \label{per:elementary-divisors-and-jordan-blocks} \hlt{Elementary Divisors and Jordan Blocks}\\
Over $\mathbb{C}$, elementary divisors are exactly $(z-\lambda)^k$, and $\mathbb{C}[z]/\langle(z-\lambda)^k\rangle$ corresponds to the Jordan block $J_k(\lambda)$.\\
Thus $\operatorname{ElemDiv}(T)$ is the multiset of Jordan blocks.
\end{perspective}

\begin{perspective} \label{per:module-and-spectral-canonical-form-crosswalk} \hlt{Crosswalk with the Spectral Presentation}\\
The module and spectral presentations encode the same data at different splitting levels:
\begin{enumerate}[label=(\roman*)]
\setlength{\itemsep}{0pt}
\item diagonalizable $\Longleftrightarrow \operatorname{ElemDiv}(T)$ consists of linear factors, by Theorem \ref{thm:diagonalizable-iff-minimal-polynomial-squarefree};
\item $(z-\lambda)^k\Longleftrightarrow J_k(\lambda)$ over an algebraically closed field, by Perspective \ref{per:elementary-divisors-and-jordan-blocks} and Theorem \ref{thm:complex-jordan-form};
\item rank drops recover the exponents $k$ via Proposition \ref{prop:p-primary-layer-counts-recover-exponents}, with $p=z-\lambda$ and $M=\mathcal{V}_T$.
\end{enumerate}
\hlt{Key Insight}: Rational form is Jordan form before irreducibles split.
\end{perspective}

\begin{perspective} \label{per:operator-module-similarity-dictionary} \hlt{Operator Modules and Similarity}\\
For $T\in\LL(\mathcal{V})$, $\mathcal{V}_T$ has action $p\cdot\mathbf{v}=p(T)\mathbf{v}$. Then $S\sim T\Longleftrightarrow\mathcal{V}_S\cong_{\mathbb{F}[z]}\mathcal{V}_T$; $\mathbb{F}[z]/\langle f\rangle\Longleftrightarrow C_f$ (Theorem \ref{thm:cyclic-subspace-quotient-and-companion}).
\end{perspective}

\begin{theorem} \label{thm:elementary-divisors-complete-similarity-invariant} \hlt{Complete Similarity Invariants}\\
Let $T,S\in\LL(\mathcal{V})$ and $A,B\in\operatorname{Mat}(n,n,\mathbb{F})$.
\begin{enumerate}[label=(\roman*)]
\setlength{\itemsep}{0pt}
\item $T\sim S\Longleftrightarrow\mathcal{V}_T\cong_{\mathbb{F}[z]}\mathcal{V}_S\Longleftrightarrow\operatorname{ElemDiv}(T)=\operatorname{ElemDiv}(S)\Longleftrightarrow\operatorname{InvFact}(T)=\operatorname{InvFact}(S)$.
\item $A\sim B\Longleftrightarrow(\mathbb{F}^n)_A\cong_{\mathbb{F}[z]}(\mathbb{F}^n)_B\Longleftrightarrow\operatorname{ElemDiv}(A)=\operatorname{ElemDiv}(B)\Longleftrightarrow\operatorname{InvFact}(A)=\operatorname{InvFact}(B)$.
\item If $A=[T]_{\mathcal{B}}$ for some basis $\mathcal{B}$ of $\mathcal{V}$, then $\operatorname{ElemDiv}(A)=\operatorname{ElemDiv}(T)$ and $\operatorname{InvFact}(A)=\operatorname{InvFact}(T)$.
\end{enumerate}
\begin{proof}
\begin{enumerate}[label=(\roman*).]
\setlength{\itemsep}{0pt}
\item[$(i)$]($\Rightarrow$) If $T\sim S$, then $\mathcal{V}_T\cong_{\mathbb{F}[z]}\mathcal{V}_S$ by Theorem \ref{thm:operator-similarity-iff-module-isomorphism}.\\
Module isomorphism preserves the PID decomposition, hence preserves $\operatorname{ElemDiv}$ and $\operatorname{InvFact}$.
\item[$(i)$]($\Leftarrow$) Equality of either $\operatorname{ElemDiv}$ or $\operatorname{InvFact}$ gives $\mathcal{V}_T\cong_{\mathbb{F}[z]}\mathcal{V}_S$ by uniqueness of the PID decomposition.\\
Thus $T\sim S$ by Theorem \ref{thm:operator-similarity-iff-module-isomorphism}.
\item[$(ii)$]($\Rightarrow$) If $A\sim B$, then $(\mathbb{F}^n)_A\cong_{\mathbb{F}[z]}(\mathbb{F}^n)_B$ by Corollary \ref{cor:matrix-similarity-iff-matrix-modules-isomorphic}; hence the divisor data agree.
\item[$(ii)$]($\Leftarrow$) Equality of either matrix divisor datum gives $(\mathbb{F}^n)_A\cong_{\mathbb{F}[z]}(\mathbb{F}^n)_B$, so $A\sim B$ by Corollary \ref{cor:matrix-similarity-iff-matrix-modules-isomorphic}.
\item[$(iii)$] If $A=[T]_{\mathcal{B}}$, then $\Phi(\mathbf{v})=[\mathbf{v}]_{\mathcal{B}}$ satisfies $\Phi(p(T)\mathbf{v})=p(A)\Phi(\mathbf{v})$.\\
Hence $\Phi:\mathcal{V}_T\to(\mathbb{F}^n)_A$ is an $\mathbb{F}[z]$-isomorphism, so the divisor data agree.
\end{enumerate}
\hlt{Structural Commentary}: Similarity is module isomorphism plus PID data.\\
\hlt{Key Insight}: Elementary divisors classify operator modules.
\end{proof}
\end{theorem}

\begin{lemma} \label{lem:determinant-of-companion-matrix} \hlt{Determinant of a Companion Block}\\
Let $f(z)=z^n+\sum_{i=0}^{n-1}c_iz^i$ be monic. Then $\det(zI_n-C_f)=f(z)$.
\begin{proof} Set $D_n(c_0,\ldots,c_{n-1})\defeq\det(zI_n-C_f)$. For $n=1$, $D_1(c_0)=z+c_0$.\\
For $n\geq2$, first-row expansion gives $D_n(c_0,\ldots,c_{n-1})=zD_{n-1}(c_1,\ldots,c_{n-1})+c_0$, since the $(1,1)$-minor is again companion-shaped and the last first-row cofactor contributes $c_0$. Induction gives $D_n=z^n+\sum_{i=0}^{n-1}c_iz^i=f(z)$.\\
\hlt{Structural Commentary}: The last column of $C_f$ stores the coefficients of the polynomial relation.\\
\hlt{Key Insight}: The shifted determinant reads that stored relation back as $f$.
\end{proof}
\end{lemma}

\begin{definition} \label{def:rational-canonical-block-forms} \hlt{Rational Canonical Block Forms}
\begin{enumerate}[label=(\roman*)]
\setlength{\itemsep}{0pt}
\item \hlt{Invariant-factor RCF}: $A=\operatorname{diag}(C_{s_1},\ldots,C_{s_r})$, where $s_j\in\mathbb{F}[z]$ are monic nonconstant and $s_r\mid\cdots\mid s_1$.
\item \hlt{Elementary-divisor RCF}: $A=\operatorname{diag}(C_{q_1},\ldots,C_{q_m})$, where each $q_j$ is power of monic irreducible polynomial.
\end{enumerate}
\end{definition}

\begin{theorem} \label{thm:rational-canonical-form-skeleton} \hlt{Rational Canonical Form Skeleton}\\
For $A\in\operatorname{Mat}(n,n,\mathbb{F})$, there is a unique tuple $\operatorname{InvFact}(A)=(s_1,\ldots,s_r)$ with $s_j\in\mathbb{F}[z]$ monic nonconstant, $s_r\mid\cdots\mid s_1$, and $A\sim_{\mathbb{F}}R_A\defeq\operatorname{diag}(C_{s_1},\ldots,C_{s_r})$. Moreover, $\mu_A=s_1$ and $p_A=\prod_{j=1}^r s_j$.\\
Equivalently, for $A,B\in\operatorname{Mat}(n,n,\mathbb{F})$, $A\sim_{\mathbb{F}}B\Longleftrightarrow\operatorname{ElemDiv}(A)=\operatorname{ElemDiv}(B)$ as multisets.
\begin{proof} Let $T_A:\mathbb{F}^n\to\mathbb{F}^n$, $\mathbf{x}\mapsto A\mathbf{x}$.\\
Theorem \ref{thm:operator-invariant-factor-skeleton} gives $\mathbb{F}^n=\bigoplus_{j=1}^r\mathcal{C}_j$, $\mathcal{C}_j\cong_{\mathbb{F}[z]}\mathbb{F}[z]/\langle s_j\rangle$, $s_r\mid\cdots\mid s_1$, $[T_A]=R_A$, $\mu_{T_A}=s_1$, and $p_{T_A}=\prod_{j=1}^r s_j$.\\
Hence $A\sim_{\mathbb{F}}R_A$, $\mu_A=s_1$, and $p_A=\prod_{j=1}^r s_j$; uniqueness of $\operatorname{InvFact}(A)$ is uniqueness in the same theorem.
\begin{enumerate}[label=\roman*.]
\setlength{\itemsep}{0pt}
\item[$(\Rightarrow)$]
If $A\sim_{\mathbb{F}}B$, then Theorem \ref{thm:elementary-divisors-complete-similarity-invariant} gives $\operatorname{ElemDiv}(A)=\operatorname{ElemDiv}(B)$.
\item[$(\Leftarrow)$]
If $\operatorname{ElemDiv}(A)=\operatorname{ElemDiv}(B)$, then Theorem \ref{thm:elementary-divisors-complete-similarity-invariant} gives $A\sim_{\mathbb{F}}B$.
\end{enumerate}
\hlt{Structural Commentary}: Companion blocks classify similarity without requiring $p_A$ to split.\\
\hlt{Key Insight}: Similarity classes can be represented by polynomial data rather than by choosing roots.
\end{proof}
\end{theorem}

\begin{theorem} \label{thm:rational-canonical-form-elementary-and-invariant-versions} \hlt{Rational Canonical Form: Two Versions}\\
Let $T\in\LL(\mathcal{V})$.
\begin{enumerate}[label=(\roman*)]
\setlength{\itemsep}{0pt}
\item If $\operatorname{ElemDiv}(T)=\{q_1,\ldots,q_m\}$, then $\exists$ basis $\mathcal{R}$ with $[T]_{\mathcal{R}}=\operatorname{diag}(C_{q_1},\ldots,C_{q_m})$.\\
The form is unique up to block permutation.
\item If $\operatorname{InvFact}(T)=(s_1,\ldots,s_r)$ and $s_r\mid\cdots\mid s_1$, then $\exists$ basis $\mathcal{B}$ with $[T]_{\mathcal{B}}=\operatorname{diag}(C_{s_1},\ldots,C_{s_r})$.\\
The form is unique in this order.
\item If $A\in\operatorname{Mat}(n,n,\mathbb{F})$, then the similarity class $[A]_{\sim_{\mathbb{F}}}$ contains both forms.\\
Elementary-divisor form is unique up to block order, invariant-factor form by its divisibility chain.
\end{enumerate}
\begin{proof}
\begin{enumerate}[label=(\roman*)]
\setlength{\itemsep}{0pt}
\item Theorem \ref{thm:operator-primary-cyclic-decomposition} gives $\mathcal{V}=\bigoplus_{\alpha=1}^m\mathcal{C}_\alpha$ with orders $q_\alpha$; Theorem \ref{thm:cyclic-subspace-quotient-and-companion} gives cyclic bases and blocks $C_{q_\alpha}$.\\
Permuting summands conjugates by a permutation matrix. 
\item Theorem \ref{thm:operator-invariant-factor-skeleton} gives $\mathcal{V}=\bigoplus_{j=1}^r\mathcal{C}_j$, $\mathcal{C}_j\cong_{\mathbb{F}[z]}\mathbb{F}[z]/\langle s_j\rangle$, and cyclic bases with block matrix $\operatorname{diag}(C_{s_1},\ldots,C_{s_r})$.
\item For $A$, apply (i),(ii) to $T_A:\mathbf{x}\mapsto A\mathbf{x}$; any chosen basis gives a similarity matrix.\\
Uniqueness follows from Theorem \ref{thm:elementary-divisors-complete-similarity-invariant}.
\end{enumerate}
\hlt{Structural Commentary}: Elementary divisors split primary blocks; invariant factors recombine them.\\
\hlt{Key Insight}: Both forms classify the same similarity classes, but they expose different polynomial data.
\end{proof}
\end{theorem}

\begin{method} \label{meth:construct-operator-from-elementary-divisors} \hlt{Constructing an Operator from Elementary Divisors}\\
Let $Q=\{q_1,\ldots,q_m\}$ be a finite multiset of monic prime powers in $\mathbb{F}[z]$, and set $n\defeq\sum_{j=1}^m\deg q_j$.
\begin{enumerate}[label=(\roman*)]
\setlength{\itemsep}{0pt}
\item Set $R_Q\defeq\operatorname{diag}(C_{q_1},\ldots,C_{q_m})\in\operatorname{Mat}(n,n,\mathbb{F})$.
\item Then $\operatorname{ElemDiv}(T_{R_Q})=Q$; hence every such $Q$ is realised on $\mathbb{F}^n$.
\item Apply Method \ref{meth:elementary-divisors-to-invariant-factors} to recombine $Q$ into $\operatorname{InvFact}(T_{R_Q})=(s_1,\ldots,s_r)$, then $\operatorname{diag}(C_{s_1},\ldots,C_{s_r})\sim_{\mathbb{F}}R_Q$.
\end{enumerate}
\hlt{Structural Commentary}: RCF constructs operators from divisor data.\\
\hlt{Key Insight}: Companion blocks realise cyclic summands; block sums realise direct sums.
\end{method}

\begin{example} \label{ex:rational-canonical-forms-from-minimal-polynomial} \hlt{Rational Canonical Possibilities from Degree Counts}\\
Let $T\in\LL(\mathbb{R}^7)$ and $\mu_T=(z-1)(z^2+1)^2$. For $Q=\operatorname{ElemDiv}(T)$, $\operatorname{lcm}Q=\mu_T$ and $\sum_{q\in Q}\deg q=7$.\\
Hence $Q_1=\{z-1,(z^2+1)^2,z^2+1\}$ or $Q_2=\{z-1,z-1,z-1,(z^2+1)^2\}$.\\
Thus the elementary-divisor RCF is $R_i=\bigoplus_{q\in Q_i}C_q$ for $i=1,2$, up to block order.
\end{example}

\begin{theorem} \label{thm:determinant-form-characteristic-polynomial} \hlt{Determinant Form of the Characteristic Polynomial}\\
For $A\in\operatorname{Mat}(n,n,\mathbb{F})$, $\chi_A(z)=\det(zI_n-A)$. If $A=[T]_{\mathcal{B}}$, then $\chi_T=\chi_A$.
\begin{proof}
Let $Q\defeq\operatorname{ElemDiv}(A)$ and choose elementary-divisor RCF $R=\bigoplus_{q\in Q}C_q\sim_{\mathbb{F}}A$.\\
By Lemma \ref{lem:determinant-of-companion-matrix}, $\det(zI-R)=\prod_{q\in Q}\det(zI-C_q)=\prod_{q\in Q}q=\chi_R$.\\
Since similarity preserves both $Q$ and $\det(zI-\cdot)$, $\chi_A=\chi_R=\det(zI-R)=\det(zI-A)$.\\
If $A=[T]_{\mathcal{B}}$, then $\operatorname{ElemDiv}(A)=\operatorname{ElemDiv}(T)$, hence $\chi_A=\chi_T$.\\
\hlt{Structural Commentary}: Rational canonical form reduces the determinant computation to companion blocks.\\
\hlt{Key Insight}: The abstract product of elementary divisors is the concrete shifted determinant.
\end{proof}
\end{theorem}

\begin{theorem} \label{thm:base-field-change-elementary-divisors} \hlt{Changing the Base Field}\\
Let $\mathbb{F}\subseteq\mathbb{K}$, $A\in\operatorname{Mat}(n,n,\mathbb{F})$, and $\operatorname{ElemDiv}_{\mathbb{F}}(A)=\{p_i^{e_{i,j}}\}_{i,j}$ with $p_i\in\mathbb{F}[z]$ monic irreducible.\\
If $p_i=\prod_{\ell=1}^{m_i}a_{i,\ell}^{d_{i,\ell}}$ in $\mathbb{K}[z]$, with $a_{i,\ell}$ pairwise distinct monic irreducible over $\mathbb{K}$, then $\operatorname{ElemDiv}_{\mathbb{K}}(A)=\{a_{i,\ell}^{d_{i,\ell}e_{i,j}}\}_{i,j,\ell}$.\\
Moreover, $\operatorname{InvFact}_{\mathbb{K}}(A)=\operatorname{InvFact}_{\mathbb{F}}(A)$.
\begin{proof}
Take elementary-divisor RCF over $\mathbb{F}$: $R=\bigoplus_{i,j}C_{p_i^{e_{i,j}}}\sim_{\mathbb{F}}A$.\\
Since $R\in\operatorname{Mat}(n,n,\mathbb{K})$ and $p_i^{e_{i,j}}=\prod_{\ell=1}^{m_i}a_{i,\ell}^{d_{i,\ell}e_{i,j}}$ with coprime factors in $\mathbb{K}[z]$,\\
Theorem \ref{thm:cyclic-composition-coprime-orders} gives $C_{p_i^{e_{i,j}}}\sim_{\mathbb{K}}\bigoplus_{\ell=1}^{m_i}C_{a_{i,\ell}^{d_{i,\ell}e_{i,j}}}$. Summing over $(i,j)$ yields $\operatorname{ElemDiv}_{\mathbb{K}}(A)$.\\
Recombination into the divisibility chain gives the original invariant factors.\\
\hlt{Structural Commentary}: Field extension refines elementary divisors.\\
\hlt{Key Insight}: Elementary divisors split; invariant factors persist.
\end{proof}
\end{theorem}

\begin{corollary} \label{cor:nonderogatory-iff-companion-similar} \hlt{Nonderogatory Means One Companion Block}\\
For $A\in\operatorname{Mat}(n,n,\mathbb{F})$, the following are equivalent:
\begin{enumerate}[label=(\roman*)]
\setlength{\itemsep}{0pt}
\item $A$ is nonderogatory, i.e. $\mu_A=\chi_A$;
\item $(\mathbb{F}^n)_A$ is cyclic;
\item $A\sim_{\mathbb{F}}C_{\chi_A}$.
\end{enumerate}
\begin{proof}
\begin{enumerate}[label=(\roman*).]
\setlength{\itemsep}{0pt}
\item[$(i)$]$\Rightarrow (ii)$ This is Theorem \ref{thm:nonderogatory-cyclic-equivalences}.
\item[$(ii)$]$\Rightarrow (iii)$ Cyclicity gives $\operatorname{InvFact}(A)=(s)$. By Theorem \ref{thm:operator-invariant-factor-skeleton}, $s=\mu_A=\chi_A$; hence $A\sim_{\mathbb{F}}C_{\chi_A}$.
\item[$(iii)$]$\Rightarrow (i)$ Since $\mu_{C_f}=f$ by Theorem \ref{thm:companion-matrix-has-prescribed-minimal-polynomial}, similarity gives $\mu_A=\mu_{C_{\chi_A}}=\chi_A$.
\end{enumerate}
\hlt{Structural Commentary}: Nonderogatory means one cyclic block.\\
\hlt{Key Insight}: One cyclic generator gives one companion block.
\end{proof}
\end{corollary}

\begin{corollary} \label{cor:diagonal-rational-canonical-form-criterion} \hlt{When the Elementary-Divisor Form Is Diagonal}\\
Let $T\in\LL(\mathcal{V})$, $Q=\operatorname{ElemDiv}(T)$, and $R_{\operatorname{ed}}=\bigoplus_{q\in Q}C_q$.\\
Then $R_{\operatorname{ed}}$ is diagonal iff $\deg q=1$ for all $q\in Q$, iff $\mu_T=\prod_{\lambda\in\Lambda}(z-\lambda)$ for some finite $\Lambda\subseteq\mathbb{F}$.
\begin{proof}
\begin{enumerate}[label=\roman*.]
\setlength{\itemsep}{0pt}
\item[$(\Rightarrow)$] If $R_{\operatorname{ed}}$ is diagonal, each block $C_q$ has size $1$, hence $\deg q=1$ for all $q\in Q$.\\
Thus $\mu_T=\operatorname{lcm}Q$ is a product of distinct linear factors.
\item[$(\Leftarrow)$] If $\mu_T=\prod_{\lambda\in\Lambda}(z-\lambda)$, then each prime power $q\in Q$ divides $\mu_T$, so $q=z-\lambda$ for some $\lambda\in\Lambda$.\\
Hence each $C_q$ is $1\times1$, and $R_{\operatorname{ed}}$ is diagonal.
\end{enumerate}
\hlt{Structural Commentary}: Diagonal RCF means all companion blocks are $1\times1$.\\
\hlt{Key Insight}: Nonlinear or repeated factors obstruct diagonal RCF.
\end{proof}
\end{corollary}

\begin{proposition} \label{prop:ab-ba-characteristic-polynomial-and-similarity} \hlt{$AB$ and $BA$ Have the Same Characteristic Polynomial}\\
For $A,B\in\operatorname{Mat}(n,n,\mathbb{F})$, $\chi_{AB}=\chi_{BA}$.\\
If $A\in\operatorname{GL}(n,\mathbb{F})$ or $B\in\operatorname{GL}(n,\mathbb{F})$, then $AB\sim_{\mathbb{F}}BA$; otherwise not necessarily.
\begin{proof}
In $\mathbb{F}(z)$, by Sylvester's identitym $\det(zI-AB)=z^n\det(I-z^{-1}AB)=z^n\det(I-z^{-1}BA)=\det(zI-BA)$, hence $\chi_{AB}=\chi_{BA}$. If $A$ is invertible, $BA=A^{-1}(AB)A$; if $B$ is invertible, $AB=B^{-1}(BA)B$.\\
Without invertibility, take $A=\begin{bmatrix}1&0\\0&0\end{bmatrix}$ and $B=\begin{bmatrix}0&1\\0&0\end{bmatrix}$.\\
Then $AB=B$, $BA=0$, $\chi_{AB}=\chi_{BA}=z^2$, but $\operatorname{rank}(AB)=1\neq0=\operatorname{rank}(BA)$, so $AB\not\sim_{\mathbb{F}}BA$.\\
\hlt{Structural Commentary}: Characteristic polynomials see less than similarity.\\
\hlt{Key Insight}: $AB$ and $BA$ share determinant data; invertibility yields similarity.
\end{proof}
\end{proposition}

\begin{example} \label{ex:base-field-rational-canonical-form-minimal-polynomial} \hlt{The Same Polynomial over Different Fields}\\
Let $\dim V=4$ and $\mu_T=(z^2+1)(z^2-2)$. Since $\deg\mu_T=4$, $T$ is nonderogatory and $\operatorname{InvFact}_K(T)=(\mu_T)$ for $K=\mathbb{Q},\mathbb{R},\mathbb{C}$; hence $R_{\operatorname{inv},K}=C_{\mu_T}$. Meanwhile
\begin{center}
\begin{tabular}{c|c}
$K$ & $\operatorname{ElemDiv}_K(T)$\\
\hline
$\mathbb{Q}$ & $\{z^2+1,\ z^2-2\}$\\
$\mathbb{R}$ & $\{z^2+1,\ z-\sqrt{2},\ z+\sqrt{2}\}$\\
$\mathbb{C}$ & $\{z-i,\ z+i,\ z-\sqrt{2},\ z+\sqrt{2}\}$
\end{tabular}
\end{center}
Thus $R_{\operatorname{ed},K}=\bigoplus_{q\in\operatorname{ElemDiv}_K(T)}C_q$: scalar extension refines elementary divisors while preserving the invariant factor. Over $\mathbb{C}$, $R_{\operatorname{ed},\mathbb{C}}$ is diagonal.
\end{example}

\begin{figure}[H]
  \centering
  \begin{tikzpicture}[
  >=Stealth,
  every node/.style={font=\small},
  class/.style={draw=blue!60!black, fill=blue!7, rounded corners=2pt,
    minimum width=2.25cm, minimum height=0.72cm, align=center, inner sep=5pt},
  data/.style={draw=dartmouthgreen!65!black, fill=dartmouthgreen!8, rounded corners=2pt,
    minimum width=2.9cm, minimum height=0.72cm, align=center, inner sep=5pt},
  form/.style={draw=gray!60, fill=gray!6, rounded corners=2pt,
    minimum width=3.45cm, minimum height=0.72cm, align=center, inner sep=5pt},
  link/.style={->, line width=0.75pt, draw=black!70},
  note/.style={font=\scriptsize, inner sep=2pt, fill=white}
]
  \node[class] (ops) at (0,0) {$[T]_{\sim}$};
  \node[class] (mods) at (4.3,0) {$[\mathcal{V}_T]_{\mathbb{F}[z]}$};
  \node[data] (invs) at (8.7,0) {$\{s_i\}:s_r\mid\cdots\mid s_1$\\$\{q_j\}$};
  \node[form] (mats) at (13.4,0) {$\operatorname{diag}(C_{s_1},\ldots,C_{s_r})$};

  \draw[link] (ops) -- node[note, above] {$p\cdot\mathbf{v}=p(T)\mathbf{v}$} (mods);
  \draw[link] (mods) -- node[note, above] {$\bigoplus_i\mathbb{F}[z]/\langle s_i\rangle$} (invs);
  \draw[link] (invs) -- node[note, above] {$s_i\mapsto C_{s_i}$} (mats);
\end{tikzpicture}

  \caption{Rational-canonical dictionary: $[T]_{\sim}$, $[\mathcal{V}_T]_{\mathbb{F}[z]}$, $\{s_i\}$ or $\{q_j\}$, and $\operatorname{diag}(C_{s_i})$ are equivalent data.}
  \label{fig:operator-module-canonical-form-dictionary}
\end{figure}
